%MIT OpenCourseWare: https://ocw.mit.edu
%RES.18-012 Algebra II Student Notes, Spring 2022
%License: Creative Commons BY-NC-SA 
%For information about citing these materials or our Terms of Use, visit: https://ocw.mit.edu/terms.

% \rhead{\rightmark}
\lhead{Lecture 31: Applications of the Galois Correspondence}

\section{Applications of the Galois Correspondence}

\subsection{Review}

Last class, we saw that if $E/F$ is a Galois extension and $G = \operatorname{Gal}(E/F)$, then there is a correspondence between subgroups $H \subset G$ and their fixed fields $E^H \subset E$. We saw that in the tower of extensions $E/E^H/F$, the top extension $E/E^H$ is always Galois, with Galois group $H$. Meanwhile, $E^H/F$ is not always Galois; but it's Galois if and only if $H$ is normal, and in that case $G/H = \operatorname{Gal}(E^H/F)$ (so in some sense, the left-hand side  makes sense if and only if the right-hand side does):

\begin{proposition}
If $K = E^H$, then $K/F$ is Galois if and only if $K$ is invariant under all $g \in G$, which occurs if and only if $H$ is normal. 
\end{proposition}

\begin{question}
    What does it mean that $K$ is invariant under all $g \in G$?
\end{question}

\begin{ans}
This means that for any $g \in G$, we have $x \in K$ if and only if $g(x) \in K$. In other words, $g(K) = K$. (So each $g$ permutes the elements of $K$; this doesn't mean that $g$ fixes each element of $K$.)
\end{ans}

\begin{question}
    Did we prove the second equivalence (that $K$ is invariant if and only if $H$ is normal)?
\end{question}

\begin{ans}
At the end of last class --- it follows from the correspondence being natural, and therefore compatible with the action of $G$. More precisely, if $H$ corresponds to $K$, then $gHg^{-1}$ corresponds to $g(K)$ (the action by $g$ on subfields corresponds to the action by $g$ on subgroups via conjugation --- this is unsurprising, since conjugation is the natural action by group elements on subgroups). From this, we see that $g(K) = K$ if and only if $gHg^{-1} = H$. 

Then $ghg^{-1}$ fixes $g(x)$ if and only if $h$ fixes $x$ --- checking this is easy, as $ghg^{-1}(g(x)) = gh(x)$. 
\end{ans}

\subsection{Cyclotomic Extensions}

The main theorem can be used to answer our question about ruler and compass constructions:

\begin{proposition}
    If $p = 2^k + 1$ is a Fermat prime, then a regular $p$-gon can be constructed by a compass and straightedge. 
\end{proposition}

\begin{proof}
    Let $\zeta$ be a $p$th root of unity. Then it suffices to show that $\QQ(\zeta)$ can be obtained by iterating quadratic extensions --- if we let $E = \QQ(\zeta)$, then it suffices to show there exists a tower of subfields \[\QQ = F_0 \subset F_1 \subset F_2 \subset \cdots \subset F_n = E,\] such that $[F_i : F_{i - 1}] = 2$ for all $i$. Quadratic extensions can always be obtained by extracting the square root of some element; so this would mean we can obtain $\QQ(\zeta)$ by starting with $\QQ$ and successively applying arithmetic operations and square roots.
    
    This is fairly clear from the Galois correspondence. We saw earlier that \[\operatorname{Gal}(E/\QQ) = (\ZZ/p\ZZ)^\times \cong \ZZ/(p - 1)\ZZ = \ZZ/2^k\ZZ.\] We can now write \[G = G_0 \supset G_1 \supset G_2 \supset \cdots \supset \{0\},\] where $G_1 = 2\ZZ/2^k\ZZ$, $G_2 = 4\ZZ/2^k\ZZ$, and so on. Then $G_i/G_{i + 1} \cong \ZZ/2\ZZ$ for all $i$. 
    
    We can then take $F_i$ to be the fixed field of $G_i$. We saw that the correspondence reverses inclusion, and we know how degrees correspond --- we have $[E : F_i] = 2^i$ for each $i$, which implies that $[F_i : F_{i - 1}] = 2$, as desired. 
\end{proof}

\begin{example}
    Describe the first step in this construction (to find $F_1$). 
\end{example}

\begin{proof}[Solution]
    We want to write down a quadratic extension of $\QQ$. We know $F_1$ is the fixed field of $G_1$, and $G_1$ consists of the even residues in the language of $\ZZ/2^k\ZZ$; converting back to the language of $(\ZZ/p\ZZ)^\times$, then $G_1$ consists of the squares (or \emph{quadratic residues}) in $\ZZ/p\ZZ$ --- elements of the form $a = b^2$ for some $b \neq 0$.
    
    Suppose $\zeta = \exp(2\pi i/p)$, and let \[\alpha = \sum_{a \in \text{QR}} \zeta^a\] (summing over all $(p - 1)/2$ quadratic residues mod $p$ --- for example, if $p = 5$, then $\alpha = \zeta + \zeta^4$). It's clear that $\alpha$ is fixed by $G_1$, since multiplying all $a$ by a quadratic residue only permutes them. 
    
    We also want to find its Galois conjugate $\beta$. To do that, we apply an element of the Galois group \emph{not} in $G_1$, which gives \[\beta = \sum_{b \in \text{NQR}} \zeta^b\] (summing over all quadratic nonresidues mod $p$ --- for example, if $p = 5$, then $\alpha = \zeta^2 + \zeta^3$.) 
    We now want to compute the quadratic equation that $\alpha$ satisfies. We know \[\alpha + \beta = \zeta^1 + \zeta^2 + \cdots + \zeta^{p - 1} = -1.\] On the other hand, we can compute \[\alpha \beta = \sum n_c\zeta^c,\] where $n_c$ is the number of ways to write $a + b = c$ where $a$ is a quadratic residue, and $b$ is a quadratic nonresidue. 
    
    This is a combinatorial problem, which we can solve --- first, $n_0 = 0$, since $-1$ is a square (this means if $a$ is a square, so is $-a$, so we can't have $a + b = 0$ where $a$ is square and $b$ isn't). On the other hand, we claim that $n_1$, \ldots, $n_{p - 1}$ are all equal --- for any $c$ and $c'$, we can write $c' = tc$ for some $t$ (since $\ZZ/p\ZZ$ is a field). If $t$ is a square, then we can get a bijection between $(a, b)$ with sum $c$ and sum $c'$, by multiplying by $t$. Meanwhile, if $t$ is not a square, then we can get a bijection by multiplying and swapping --- given $(a, b)$ with sum $c$, we can take $(tb, ta)$ with sum $c'$. This means $n_c = n_{c'}$. Finally, we have $n_0 + \cdots + n_{p - 1} = ((p - 1)/2)^2$, since this is the number of ways to choose a summand from each of $\alpha$ and $\beta$. This means \[n_c = \begin{cases} \frac{p - 1}{4} & \text{if $c \neq 0$} \\ 0 & \text{if $c = 0$},\end{cases}\] so then our sum is \[\alpha\beta = \sum_{c = 1}^{p - 1} \frac{p - 1}{4}\zeta^c = -\frac{p - 1}{4}.\] This means our quadratic equation is \[\alpha^2 + \alpha - \frac{p - 1}{4} = 0 \implies \alpha = \frac{-1 \pm \sqrt{p}}{2}.\] So we have $F_1 = \QQ(\sqrt{p})$.  
\end{proof}

\begin{note}
    This argument works for any prime $p \equiv 1 \pmod{4}$, meaning that the quadratic extension of $\QQ$ contained in $\QQ(\zeta_p)$ is still $\QQ(\sqrt{p})$. Meanwhile, when $p \equiv 3 \pmod{4}$, we instead get $\QQ(\sqrt{-p})$. 
\end{note}

The description of $\operatorname{Gal}(\QQ(\zeta)/\QQ)$ can be generalized to apply to all $n$ (meaning $\zeta$ is an $n$th root of unity), not just primes. 

\begin{definition}
    The $n$th \textbf{cyclotomic polynomial} $\Phi_n$ is the monic polynomial in $\mathbb{Z}[x]$ whose roots are exactly the primitive $n$th roots of unity.
\end{definition}

We then have \[x^n - 1 = \prod_{d \mid n} \Phi_d(x).\] This is because the roots of $x^n - 1$ are all elements whose order in $\CC^\times$ divides $n$, and the right-hand side groups such terms by their order $d$). 

This formula lets us compute $\Phi_n$. 

\begin{example}
    We have $\Phi_1(x) = x - 1$, and \[\Phi_p(x) = x^{p - 1} + \cdots + 1.\] We can also compute other polynomials $\Phi_n(x)$, such as \[\Phi_{12}(x) = x^4 - x^2 + 1.\] 
\end{example}

The cyclotomic polynomials don't always have all coefficients $0$ or $\pm 1$, but the smallest counterexample is $105$ (the smallest product of three distinct odd primes). But from this formula, it's easy to show by induction that all $\Phi_n$ have integer coefficients. 

\begin{fact}
    $\Phi_n$ is irreducible in $\mathbb{Q}[x]$. 
\end{fact}

We proved this fact for primes; we won't prove it for general $n$, since the proof is longer. 

% START HERE

Also note that $\deg(\Phi_n)$ is the number of elements of order $n$ in the additive group $\ZZ/n\ZZ$, which is $\varphi(n) = \lvert (\ZZ/n\ZZ)^\times \rvert$. If $n = p_1^{d_1}\cdots p_k^{d_k}$, we have the explicit formula \[\varphi(n) = \prod_i (p_i^{d_i} - p_i^{d_i - 1}).\] 

Now we have \[[\QQ(\zeta) : \QQ] = \varphi(n),\] and $\QQ(\zeta)$ is a splitting field (for the same reason as in the prime case --- all roots of $\Phi_n$ are powers of $\zeta$). By the same reasoning as the prime case, we then have \[\operatorname{Gal}(\QQ(\zeta)/\QQ) = (\ZZ/n\ZZ)^*.\] Note that this is not necessarily cyclic --- in fact, it's not cyclic unless $n$ is a prime power or twice a prime power (and it's also not cyclic if $n \geq 8$ is a power of $2$). It'll be the \emph{product} of cyclic groups (since it's still abelian), but there will usually be multiple factors of even order in this product. 

\subsection{Kummer Extensions}

We'll now consider extensions $E/F$ where $E = F(\alpha)$ for some $\alpha$ such that $\alpha^n \in F$ for a positive integer $n$ (and $\alpha \neq 0$). Assume that $F$ contains all $n$th roots of unity, meaning that \[\mu_n(F) = \{x \in F \mid x^n = 1\}\] has exactly $n$ elements (and therefore $\mu_n(F) \cong \ZZ/n\ZZ$); this is equivalent to requiring that $F$ contains a primitive $n$th root of $1$. 

Our main example is over characteristic $0$, but this can be done over characteristic $p$ as well, with the additional requirement that $p \nmid n$. 

\begin{proposition}
    In this case $E/F$ is Galois, and \[\operatorname{Gal}(E/F) \cong \mathbb{Z}/m\ZZ\] for some $m \mid n$. In fact, if $x^n - a$ is irreducible in $F[x]$, then $m = n$. 
\end{proposition}

\begin{proof}
    We have \[x^n - a = \prod (x - \zeta^i\alpha),\] where $0 \leq i \leq n - 1$ and $\zeta$ is a primitive $n$th root of $1$ (since all $\zeta^i\alpha$ are roots of $x^n - a$, and they are all distinct). So if we're given one root of $x^n - a$, then all possible roots are obtained by multiplication by roots of unity (which are in $F$). So $E$ is the splitting field of $x^n - a$. 
    
    Now an element $\sigma \in G = \operatorname{Gal}(E/F)$ is uniquely determined by $\sigma(\alpha)$, which must be $\zeta^i\alpha$ for some $i$. For each $i$, let $\sigma_i$ be the element in $G$ such that $\sigma_i(\alpha) = \zeta^i\alpha$, if it exists (the element $\sigma_i$ doesn't necessarily exist for all $i$). 

    It's clear that \[\sigma_i\sigma_j(\alpha) = \sigma_i(\zeta^j\alpha) = \zeta^{i + j}\alpha = \sigma_{i + j}(\alpha)\] (because $\zeta \in F$, so $\sigma$ must fix it). So then $\sigma_i\sigma_j = \sigma_{i + j}$. This means $G$ is isomorphic to a subgroup in $\mathbb{Z}/n\ZZ$, and every such subgroup must be of the form $\ZZ/m\ZZ$ where $m \mid n$.  

    In fact $m = \deg(E/F)$, so $m = n$ if and only if $x^n - a$ is irreducible. (When $x^n - a$ to be reducible, this fails in a trivial way --- then a \emph{smaller} power of $\alpha$ is in $F$.)
\end{proof}

\subsection{Quintic Equations}

Using these ideas, we can obtain the famous application of Galois theory to the impossibility of solving a general polynomial equation of degree at least $5$. 

\begin{definition}
    A finite group $G$ is \textbf{solvable} if there exists a sequence of subgroups \[G = G_0 \supset G_1 \supset G_2 \supset \cdots \supset G_n = \{1\}\] such that for all $i$, $G_i$ is a normal subgroup of $G_{i - 1}$ and $G_{i - 1}/G_i$ is abelian. 
\end{definition}

The main idea of the proof is the following two propositions: 

\begin{proposition}
    Given an extension $E/F$ and some $\alpha \in E$ such that $\alpha$ can be obtained from elements of $F$ by arithmetic operations (addition, subtraction, multiplication, and division) and extracting arbitrary $n$th roots (where we're allowed to choose any of the possible $n$th roots), then $\alpha$ lies in a Galois extension of $F$ with a solvable Galois group. 
\end{proposition}

\begin{proposition}
    $S_n$ is not solvable for $n \geq 5$. 
\end{proposition}

The first proposition essentially follows from what we've already discussed --- we'll discuss it in more detail next class, but the idea is to first add the roots of unity; then when we extract a $n$th root, we get an extension with cyclic Galois group. Then when we extract $n$th roots repeatedly, we get a sequence of subgroups with abelian quotients. Meanwhile, the second is an elementary finite group argument. 

\begin{corollary}
    A root of a polynomial $P$ of degree $5$ with Galois group $S_5$ cannot be expressed through the rational numbers in radicals. 
\end{corollary}

Saying the root can't be expressed in radicals is shorthand for the longer sentence from earlier --- it simply means that it can't be obtained by arithmetic operations and extracting $n$th roots. 

So this means not only is there no universal formula for the roots using radicals (as there is in lower degrees), there isn't even a way to write down the roots of a \emph{specific} polynomial. 

\begin{proof}[Proof of corollary]
If it were possible to express all roots of $P$ in radicals, then the splitting field $K$ of $P$ would be contained in a Galois extension of $\QQ$ with solvable Galois group $G$. But then we have an onto homomorphism $G \twoheadrightarrow \operatorname{Gal}(K/\QQ) = S_5$. But the quotient of a solvable group is again solvable; so this would imply $S_5$ is solvable, contradiction. 
\end{proof}
